\clearpage
\VerdierPhysicalPageAnchor{0044}%
\VerdierSemanticAnchor{ch1-1-4}
\subsection{群的变换}

\VerdierSemanticAnchor{ch1-1-4-1}
\noindent\textbf{1.4.1.}\quad
设 $f\colon H\to G$ 为一个群同态。对每个 $G$-范畴
$(\mathcal C,G,T,c)$，令 $H$ 通过 $f$ 作用于 $\mathcal C$，即置

\VerdierSemanticAnchor{ch1-1-4-1-eq1}
\begin{equation}
  \begin{aligned}
    T_f(h)&=T\bigl(f(h)\bigr),
      & h&\in H,\\
    c_f(h,h')&=c\bigl(f(h),f(h')\bigr),
      & h,h'&\in H.
  \end{aligned}
  \tag{1.4.1.1}
\end{equation}
便得到一个 $H$-范畴 $(\mathcal C,H,T_f,c_f)$。

\VerdierSemanticAnchor{ch1-1-4-1-gfunctor}
同样，给定一个 $G$-函子
\[
  (F,m)\colon(\mathcal C,G,T,c)
  \longrightarrow(\mathcal C',G,T',c'),
\]
置

\VerdierSemanticAnchor{ch1-1-4-1-eq2}
\begin{equation}
  m_f(h)=m\bigl(f(h)\bigr),
  \tag{1.4.1.2}
\end{equation}
便得到相应 $H$-范畴之间的一个 $H$-函子 $(F,m_f)$。

\VerdierSemanticAnchor{ch1-1-4-1-morphisms}
最后，一个与 $G$ 的作用相容的 $G$-函子态射也与 $H$ 的作用相容。

\VerdierSemanticAnchor{ch1-1-4-1-eq3}
上述构造显然定义了一个函子，称为\emph{换群函子}：
\begin{equation}
  f^*\colon G\text{-}\operatorname{Cat}
  \longrightarrow H\text{-}\operatorname{Cat}.
  \tag{1.4.1.3}
\end{equation}

\VerdierSemanticAnchor{ch1-1-4-1-eq4}
更精确地，设 $\mathcal C$ 与 $\mathcal C'$ 为两个 $G$-范畴；函子 $f^*$
定义一个映射
\begin{equation}
  \begin{aligned}
    f^*(\mathcal C,\mathcal C')\colon\quad
      &\operatorname{Hom}_{G\text{-}\operatorname{Cat}}
        (\mathcal C,\mathcal C')\\
      &\longrightarrow
        \operatorname{Hom}_{H\text{-}\operatorname{Cat}}
        \bigl(f^*(\mathcal C),f^*(\mathcal C')\bigr).
  \end{aligned}
  \tag{1.4.1.4}
\end{equation}

\VerdierSemanticAnchor{ch1-1-4-1-homcategories}
但集合\allowbreak
$\operatorname{Hom}_{G\text{-}\operatorname{Cat}}(\mathcal C,\mathcal C')$\allowbreak
与\allowbreak
$\operatorname{Hom}_{H\text{-}\operatorname{Cat}}
\bigl(f^*(\mathcal C),f^*(\mathcal C')\bigr)$\allowbreak
分别是范畴\hfill\break
$\mathcal H\!om_{G\text{-}\operatorname{Cat}}(\mathcal C,\mathcal C')$\allowbreak
与\allowbreak
$\mathcal H\!om_{H\text{-}\operatorname{Cat}}
\bigl(f^*(\mathcal C),f^*(\mathcal C')\bigr)$\allowbreak
的对象集；这些范畴的态射是与群作用相容的函子态射，而且映射
$f^*(\mathcal C,\mathcal C')$ 自然地延拓为函子

\VerdierSemanticAnchor{ch1-1-4-1-eq5}
\begin{equation}
  \begin{aligned}
    f^*(\mathcal C,\mathcal C')\colon\quad
      &\mathcal H\!om_{G\text{-}\operatorname{Cat}}
        (\mathcal C,\mathcal C')\\
      &\longrightarrow
        \mathcal H\!om_{H\text{-}\operatorname{Cat}}
        \bigl(f^*(\mathcal C),f^*(\mathcal C')\bigr).
  \end{aligned}
  \tag{1.4.1.5}
\end{equation}

\VerdierSemanticAnchor{ch1-1-4-1-faithful}
这个函子是忠实的。
